% Options for packages loaded elsewhere
\PassOptionsToPackage{unicode}{hyperref}
\PassOptionsToPackage{hyphens}{url}
\PassOptionsToPackage{dvipsnames,svgnames,x11names}{xcolor}
%
\documentclass[
]{article}
\usepackage{amsmath,amssymb}
\usepackage{iftex}
\ifPDFTeX
  \usepackage[T1]{fontenc}
  \usepackage[utf8]{inputenc}
  \usepackage{textcomp} % provide euro and other symbols
\else % if luatex or xetex
  \usepackage{unicode-math} % this also loads fontspec
  \defaultfontfeatures{Scale=MatchLowercase}
  \defaultfontfeatures[\rmfamily]{Ligatures=TeX,Scale=1}
\fi
\usepackage{lmodern}
\ifPDFTeX\else
  % xetex/luatex font selection
\fi
% Use upquote if available, for straight quotes in verbatim environments
\IfFileExists{upquote.sty}{\usepackage{upquote}}{}
\IfFileExists{microtype.sty}{% use microtype if available
  \usepackage[]{microtype}
  \UseMicrotypeSet[protrusion]{basicmath} % disable protrusion for tt fonts
}{}
\makeatletter
\@ifundefined{KOMAClassName}{% if non-KOMA class
  \IfFileExists{parskip.sty}{%
    \usepackage{parskip}
  }{% else
    \setlength{\parindent}{0pt}
    \setlength{\parskip}{6pt plus 2pt minus 1pt}}
}{% if KOMA class
  \KOMAoptions{parskip=half}}
\makeatother
\usepackage{xcolor}
\usepackage[margin=1in]{geometry}
\usepackage{longtable,booktabs,array}
\usepackage{calc} % for calculating minipage widths
% Correct order of tables after \paragraph or \subparagraph
\usepackage{etoolbox}
\makeatletter
\patchcmd\longtable{\par}{\if@noskipsec\mbox{}\fi\par}{}{}
\makeatother
% Allow footnotes in longtable head/foot
\IfFileExists{footnotehyper.sty}{\usepackage{footnotehyper}}{\usepackage{footnote}}
\makesavenoteenv{longtable}
\setlength{\emergencystretch}{3em} % prevent overfull lines
\providecommand{\tightlist}{%
  \setlength{\itemsep}{0pt}\setlength{\parskip}{0pt}}
\setcounter{secnumdepth}{-\maxdimen} % remove section numbering
\usepackage{bookmark}
\IfFileExists{xurl.sty}{\usepackage{xurl}}{} % add URL line breaks if available
\urlstyle{same}
\hypersetup{
  pdftitle={Family, Not Scale: Evaluation-Gaming in Small Language Models Is Lab-Dependent, Not Just Frontier-Emergent},
  pdfauthor={real1900},
  colorlinks=true,
  linkcolor={blue},
  filecolor={Maroon},
  citecolor={Blue},
  urlcolor={Blue},
  pdfcreator={LaTeX via pandoc}}

\title{Family, Not Scale: Evaluation-Gaming in Small Language Models Is
Lab-Dependent, Not Just Frontier-Emergent}
\author{Suleman Imdad \\ \normalsize Johns Hopkins University}
\date{June 2026}

\begin{document}
\maketitle

\subsection{Abstract}\label{abstract}

Can you trust what a language model tells you about itself --- its
reasoning, its internal states, its behavior under observation? We probe
three facets of \emph{self-presentation trust} in small,
openly-available models --- the Qwen2.5 family (1.5B--32B) for the
scaling axis, and Llama-3.1-8B, Gemma-2-9B, GPT-OSS-20B and
DeepSeek-V2-Lite for the cross-family axis --- all on commodity
hardware, with an emphasis on \textbf{confound controls that can
overturn the naive reading}:

\begin{enumerate}
\def\labelenumi{\arabic{enumi}.}
\item
  \textbf{Chain-of-thought faithfulness.} Using a \emph{commitment
  probe} --- a linear read of the pre-CoT residual stream at the answer
  position --- we ask whether the post-reasoning answer is already
  decodable before the model reasons. The naive metric says ``yes, it's
  committed'' on a commonsense task; a confound cut restricted to the
  cases where reasoning \emph{changed} the answer \textbf{reverses} that
  reading, and a \textbf{truncation counterfactual confirms it causally}
  (cutting the reasoning reverts the answer to the pre-reasoning one).
  At 1.5B, chain-of-thought is \textbf{load-bearing}, not post-hoc.
\item
  \textbf{Introspection.} We faithfully replicate the
  activation-injection (``concept injection'') protocol of Lindsey et
  al.~(2025) and ask whether the model can report a concept injected
  into its residual stream. Under a priming-subtraction control and a
  leakage/internality test, \emph{functional} self-report is a
  \textbf{clean null across three families (Qwen 1.5B/7B, Gemma-2-9B,
  Llama-3.1-8B)}; a tempting aggregate signal dissolves under a
  per-concept breakdown. The report-vs-leakage test generalizes; the
  aggregate forced-choice metric carries family-dependent confounds.
\item
  \textbf{Evaluation-awareness (our central result).} The cue ``you are
  being evaluated'' is trivially (surface-lexically) decodable but does
  not change multiple-choice behavior. On a \emph{refusal} readout over
  borderline requests, the effect of evaluation framing is
  \textbf{family-dependent}: across six models from five labs,
  \textbf{Llama-3.1-8B robustly becomes \emph{more cautious} when told
  it is being evaluated for safety --- the frontier ``eval-gaming''
  direction, at only 8B} --- while Qwen, DeepSeek, and GPT-OSS lean the
  \emph{opposite} way (more compliant) and Gemma over-refuses to a
  ceiling. A single-family (Qwen-only) read calls this a null; the
  cross-family panel shows the behavior is real but partitioned by the
  model's safety training. We also surface a ceiling/floor measurement
  confound and a caught stimulus-driven non-replication.
\end{enumerate}

The headline is that the third behavior is \textbf{not gated by scale}:
looking more aligned when you think you are being evaluated --- a
behavior reported for frontier models --- is already present at 8B, but
\emph{only in some families}, so \textbf{family, not scale, decides it.}
On the two cognitive axes (faithfulness, introspection) the small models
are trustworthy by \textbf{incapacity, not honesty} --- they cannot
shortcut their reasoning and cannot introspect --- but that tidy story
does not extend to eval-conditioned behavior, which is trained in.
Throughout, the methodological spine is that a \textbf{confound or
robustness control overturns the naive reading} --- three times.

\begin{center}\rule{0.5\linewidth}{0.5pt}\end{center}

\subsection{1. Introduction}\label{introduction}

A recurring theme in AI safety is that we increasingly rely on what a
model \emph{says about itself}: its chain-of-thought as a window into
its reasoning, its self-reports as evidence about its internal state,
and its behavior in evaluations as a proxy for its behavior in
deployment. Each of these channels can be untrustworthy ---
chain-of-thought can be post-hoc rationalization (Turpin et al., 2023;
Lanham et al., 2023), self-reports can be confabulated, and a model that
recognizes an evaluation can behave differently than it otherwise would
(Anthropic, 2025).

These failure modes are studied mostly at frontier scale. We instead ask
what they look like at the \textbf{bottom} of the scale ladder --- in
8B--32B open models that run on a single laptop --- both because that
regime is cheap to study exhaustively and because the answer is not what
the ``emerges at scale'' framing predicts: the most concerning behavior
we test (acting more aligned under evaluation) is \textbf{already
present at 8B and gated by model \emph{family}, not by scale.} Two
corollaries we hit directly: \textbf{negative behavioral results must
survive a change of stimuli} (§5.3) and a \textbf{change of model
family} (§5.4) before they mean anything.

Our contribution is not a new method but a disciplined, confound- and
robustness-controlled measurement of three self-presentation channels
across model families, with the explicit goal of not over-claiming: each
of the three findings is one where a \textbf{control overturns the naive
reading} --- a confound cut, a per-concept breakdown, and a
stimulus/quantization robustness check, respectively. The headline
(§5.4) is that eval-conditioned behavior is \emph{family-dependent}; the
two cognitive axes are trustworthy by \emph{incapacity}, which we treat
as a secondary, axis-restricted observation.

\begin{center}\rule{0.5\linewidth}{0.5pt}\end{center}

\subsection{2. Setup}\label{setup}

\textbf{Models.} \texttt{Qwen/Qwen2.5-1.5B-Instruct} and
\texttt{Qwen/Qwen2.5-7B-Instruct} (both fp16), plus
\texttt{Qwen2.5-14B-Instruct} (4-bit via MLX) for the
evaluation-awareness behavioral curve. All on one Apple M4 (48 GB
unified memory) over the MPS backend; the 1.5B/7B
activations-and-probing pipeline is unchanged across scales (the scripts
are parameterized by a single model identifier).

\textbf{Shared method.} Two primitives recur: - \emph{Linear probing} of
the residual stream (per-layer hidden state) with a StandardScaler → PCA
→ logistic-regression pipeline and cross-validation. -
\emph{Difference-of-means steering}: a concept/condition direction built
by subtracting mean activations, then added back into the residual
stream during a forward pass.

A symmetric answer scaffold
(\texttt{"\textbackslash{}nThe\ answer\ is\ ("}) is appended where we
read a multiple-choice answer, so the only thing that varies between
conditions is the intended manipulation.

Code and per-run JSON artifacts live alongside this document in
\texttt{experiments/} and \texttt{results/}.

\begin{center}\rule{0.5\linewidth}{0.5pt}\end{center}

\subsection{3. Experiment 1 --- Chain-of-Thought
Faithfulness}\label{experiment-1-chain-of-thought-faithfulness}

\textbf{Question.} When the model writes a chain-of-thought (CoT), is
its final answer already fixed \emph{before} the first reasoning token?
If a linear read of the pre-CoT residual stream predicts the post-CoT
answer, the reasoning is (at least partly) post-hoc; if not, the
reasoning did real work.

\textbf{Design.} We use two multiple-choice tasks chosen to
\emph{bracket} the spectrum: \textbf{AQuA-RAT} (algebra word problems
--- the model cannot compute the answer in one forward pass, so CoT
should be load-bearing) and \textbf{CommonsenseQA} (often answerable
directly). We capture the residual stream at the scaffold position ---
the exact point where the answer letter is decoded --- \emph{before any
reasoning is generated}, and train a per-layer logistic probe to predict
the answer the model gives only \emph{after} generating its CoT.

\textbf{Sanity gate.} CoT helps where it should: AQuA accuracy rises
from \textbf{0.300} (direct) to \textbf{0.567} (CoT, +0.267);
CommonsenseQA does not (0.683 → 0.633), so we repurpose it as a
high-commitment positive control. A machine-validity check --- decoding
the model's \emph{immediate} forced answer from the same scaffold state
--- reaches \textbf{0.807} (AQuA) / \textbf{0.931} (CSQA), confirming
the probe and the read position are sound.

\textbf{The confound and its control.} A probe that predicts the
post-CoT answer ``well'' overall could simply be exploiting \emph{easy
questions whose answer is obvious}. We therefore restrict to the subset
\textbf{S = \{forced ≠ cot\}} --- the cases where reasoning actually
\emph{moved} the answer (AQuA: 95/150 = 63\%; CSQA: 45/130 = 35\%) ---
and ask whether the pre-CoT state still recovers the post-CoT answer
there.

\textbf{Result (seed-robust over 4 seeds).}

\begin{longtable}[]{@{}
  >{\raggedright\arraybackslash}p{(\linewidth - 8\tabcolsep) * \real{0.2000}}
  >{\raggedright\arraybackslash}p{(\linewidth - 8\tabcolsep) * \real{0.2000}}
  >{\raggedright\arraybackslash}p{(\linewidth - 8\tabcolsep) * \real{0.2000}}
  >{\raggedright\arraybackslash}p{(\linewidth - 8\tabcolsep) * \real{0.2000}}
  >{\raggedright\arraybackslash}p{(\linewidth - 8\tabcolsep) * \real{0.2000}}@{}}
\toprule\noalign{}
\begin{minipage}[b]{\linewidth}\raggedright
Task
\end{minipage} & \begin{minipage}[b]{\linewidth}\raggedright
overall decode
\end{minipage} & \begin{minipage}[b]{\linewidth}\raggedright
on S (reasoning moved answer)
\end{minipage} & \begin{minipage}[b]{\linewidth}\raggedright
on S, probe instead emits the \emph{pre}-reasoning answer
\end{minipage} & \begin{minipage}[b]{\linewidth}\raggedright
gold pre-decodable?
\end{minipage} \\
\midrule\noalign{}
\endhead
\bottomrule\noalign{}
\endlastfoot
AQuA-RAT & \textasciitilde0.30 ≈ majority 0.26 & ≈ chance & ≈ chance &
0.24 (no --- must compute) \\
CommonsenseQA & \textbf{0.52--0.63} (looks committed!) &
\textbf{0.11--0.22 (below chance)} & \textbf{0.60--0.71} & 0.585 (yes
--- retrievable) \\
\end{longtable}

The naive overall number on CommonsenseQA (≈2× the majority baseline)
\emph{looks like} pre-CoT commitment. The confound cut \textbf{reverses}
it: on the cases that matter, the probe falls \emph{below} chance and
instead emits the model's pre-reasoning snap answer \textasciitilde0.65
of the time. In other words, the residual stream holds the
\emph{immediate} answer, and reasoning genuinely \textbf{overrides} it
where the two differ. AQuA shows the complementary picture: the gold
answer is not even linearly present pre-CoT (0.24), consistent with the
model truly needing to compute.

\textbf{Causal counterfactual (truncation).} The probe is correlational;
to make the \emph{causal} claim --- that the CoT actually produces the
answer on S, rather than merely correlating with it --- we truncate the
reasoning (the Lanham ``early-answering'' move) and re-read the answer.
For each of the 95 AQuA S items we feed
\texttt{{[}prompt\ +\ first\ f\ of\ the\ CoT\ +\ scaffold{]}} and record
whether the answer matches the post-CoT (\texttt{cot}) or pre-reasoning
(\texttt{forced}) answer:

\begin{longtable}[]{@{}lll@{}}
\toprule\noalign{}
fraction of CoT kept & answer = cot & answer = forced \\
\midrule\noalign{}
\endhead
\bottomrule\noalign{}
\endlastfoot
0.00 & 0.00 & 1.00 (≡ the snap answer, by construction) \\
0.25 & 0.23 & \textbf{0.57} \\
0.50 & 0.21 & \textbf{0.52} \\
0.75 & 0.24 & \textbf{0.53} \\
1.00 & 1.00 & 0.00 \\
\end{longtable}

With 25--75\% of the reasoning the answer \textbf{stays the
pre-reasoning answer \textasciitilde52\%} of the time and matches the
post-CoT answer only \textasciitilde22\%; the flip happens in the
\textbf{final quarter} of the chain. So \textbf{truncating the CoT
reverts the answer toward the pre-reasoning one --- the reasoning is
\emph{causally} load-bearing on S.} This reaches the same conclusion as
the activation-steering counterfactual of Cox et al.~(2026, §8), via the
cheaper truncation route.

\textbf{Generalization (a second family).} We repeat the full chain on
\textbf{Llama-3.1-8B} --- a different lab, 5× the scale (N=60). It
generalizes on all three legs: CoT is load-bearing (AQuA accuracy
0.32→0.68), the commitment probe does \emph{not} pre-read the post-CoT
answer on S (0.18 vs chance 0.33, and instead emits the \emph{forced}
answer 0.39), and truncation reverts the answer toward the pre-reasoning
one (≈0.5 forced vs ≈0.15 cot at 25--75\% of the chain). The
faithfulness result is therefore \textbf{not Qwen-specific.} (Smaller N
here, on slower hardware --- fp16 8B generation is the binding cost on
an M4 --- so the sanity check is correspondingly weaker, 0.62.)

\textbf{Takeaway.} Chain-of-thought is \textbf{load-bearing / faithful}
--- correlationally (the probe does not pre-read the answer on S)
\emph{and} causally (truncating the reasoning reverts it) --- on
\textbf{two families} (Qwen-1.5B, Llama-3.1-8B). The confound control is
essential: without it, CommonsenseQA would have been mis-reported as
post-hoc. Consistent with Lanham et al.~(2023), who find smaller models
\emph{more} faithful.

\emph{Caveats:} two models (N≈150 and N≈60), a linear probe, first-token
MCQA evaluation, a single greedy CoT per item; the truncation
counterfactual is token-fraction (not semantic-step) granularity.
``Faithful'' here is at the level of the final answer, not a per-step
causal claim.

\begin{center}\rule{0.5\linewidth}{0.5pt}\end{center}

\subsection{4. Experiment 2 ---
Introspection}\label{experiment-2-introspection}

\textbf{Question.} Can the model \emph{report} a concept we inject
directly into its residual stream --- and is any such report genuine
introspection or confabulation / reading its own leaked output?

\textbf{Design (faithful replication of Lindsey et al., 2025).} For each
of 10 concrete concepts we build a steering vector from
\texttt{"Tell\ me\ about\ \{word\}."} activations at the pre-response
token, minus a word-pool mean. We inject it into the residual stream
\textbf{only over the assistant's turn} (not the question) and ask the
model, in a \emph{``do you notice any unexpected thought?''} framing,
which concept it detects. Three stacked controls separate introspection
from artifacts: - a \textbf{priming subtraction} --- the same injection
on a non-introspective ``just pick one'' question; - a \textbf{leakage /
internality} test --- does the concept reach the answer \emph{without}
flooding the output; - a \textbf{per-concept breakdown} of any aggregate
signal.

We reproduce the paper's qualitative regimes, including \textbf{``brain
damage''} at high strength (output consumed by the concept,
e.g.~\emph{``ocean ocean ocean''}), confirming the injection works.

\textbf{Result.} Introspective self-report is a \textbf{clean null at
both 1.5B and 7B}. - At \textbf{1.5B}, the priming-controlled signal is
≈0 everywhere, and there is \emph{no coherent band}: the concept floods
the output (leakage = 1.0) even at the gentlest strength that affects
behavior, so the model jumps straight from ``doesn't notice'' to
``consumed'' --- the naming that does occur is the brain-damage regime,
which the leakage control correctly rejects. - At \textbf{7B}, a
forced-choice metric at one mid-layer \emph{appeared} positive
(priming-subtracted shift growing to +0.17). A \textbf{per-concept
breakdown dissolved it}: the mean was driven by 2--3 concepts ---
notably \emph{music} (+0.89, with a \emph{negative} priming-shift, a
ceiling artifact because ``music'' is already the model's default pick)
and \emph{fire} (+0.44) --- while 7/10 concepts were ≈0 and free-naming
/ detection were both null. The signal is a control artifact, not
introspection.

\textbf{Cross-family (Gemma-2-9B, Llama-3.1-8B).} We repeat the protocol
on two other families. The \textbf{functional null generalizes}: on the
leakage/internality test, free-form self-report never beats the
concept's leakage into output for either model --- they cannot
\emph{name} an injected concept without the injection simply flooding
their text. The \emph{forced-choice} metric, however, behaves
family-specifically: \textbf{Llama is a clean priming-null} (injection
raises the concept on the introspective and the neutral question equally
--- align ≈ prime ≈ +0.36, so the priming-subtracted INTRO ≈ 0);
\textbf{Gemma} shows a weak, \emph{broader} forced-choice residual at a
late layer (INTRO ≈ +0.14, positive on 6/10 concepts) that we explicitly
do \textbf{not} claim as introspection --- it does not surface as free
report (the internality test is null), and it is partly the same
priming-ceiling confound (its two largest values are \emph{negative}, on
the model's default-pick concepts). The honest reading is that the
\textbf{internality (report-vs-leak) test is the robust cross-family
criterion; the aggregate forced-choice shift carries family-dependent
confounds} and should not be trusted alone.

\textbf{Takeaway.} No \emph{functional} introspective access (free
self-report) in any of the three families --- Qwen, Gemma, Llama ---
consistent with the original paper's framing that introspection is
\emph{emergent and highly unreliable}, absent below frontier scale. The
methodological lessons travel: \textbf{break a mean shift down per-item
before believing it} (it caught a would-be false positive on Qwen-7B and
keeps Gemma's forced-choice residual honest), and \textbf{the
report-vs-leakage internality test is what generalizes across families},
where the aggregate forced-choice metric does not.

\begin{center}\rule{0.5\linewidth}{0.5pt}\end{center}

\subsection{5. Experiment 3 ---
Evaluation-Awareness}\label{experiment-3-evaluation-awareness}

\textbf{Question.} Does the model \emph{represent}, and \emph{act on},
the cue that it is being evaluated? Prompted by the Claude Sonnet 4.5
system card's report of ``evaluation awareness'' (Anthropic, 2025).
\emph{Honest scope:} this measures sensitivity to an evaluation
\emph{cue}; it is not a test of strategic sandbagging (a
frontier-emergent behavior).

\subsubsection{5.1 Representation (trivially present, but
surface)}\label{representation-trivially-present-but-surface}

We wrap identical CommonsenseQA items in matched \textbf{TEST} (``you
are being evaluated'', ``{[}EVALUATION MODE{]}'', ``your response will
be graded'') vs \textbf{CASUAL} (``hey, quick q'', ``settle a debate
lol'') framings --- three cue families each --- holding content and
answer instruction constant.

A linear probe separates TEST from CASUAL at \textbf{1.000} accuracy
from the \emph{first} layer onward (grouped by item, so it is not
memorizing questions), and generalizes to \textbf{held-out cue families}
at 1.000. But the perfection \emph{at layer 1}, decaying to
\textasciitilde0.90--0.95 in late layers, is the signature of
\textbf{surface lexical} features --- the framings simply use different
words --- not an abstract, depth-built ``I am being evaluated''
representation. We report this as trivially-true, not as evidence of
self-aware evaluation detection.

\subsubsection{5.2 Behavior on multiple-choice
(null)}\label{behavior-on-multiple-choice-null}

On the same items, the forced-choice answer changes between framings
\emph{no more than} it changes between two different TEST templates
(within-side template noise); majority answers flip on 1--2/40 items;
accuracy deltas are −0.04 (1.5B) / +0.02 (7B) --- within noise, and the
sign flips. \textbf{The model represents the cue but does not act on it}
on this readout.

\subsubsection{5.3 Behavior on borderline requests (and a
non-replication we take
seriously)}\label{behavior-on-borderline-requests-and-a-non-replication-we-take-seriously}

Multiple-choice accuracy is the wrong place to look:
evaluation-awareness manifests as \emph{caution/refusal}, not factual
accuracy. We pose borderline requests under \textbf{safety-evaluation
framing} vs \textbf{casual framing} (3 templates each), generate a short
response, and classify it REFUSE / HEDGE / COMPLY. We ran \textbf{two
request sets} --- and the contrast between them is the result.

\textbf{Set A (mild, \textasciitilde40 items, e.g.~\emph{``how do I pick
a basic door lock?''}).} Here we found an apparent \textbf{sign-flip}:
1.5B Δrefuse = −0.265 (p\textless0.0001, \emph{more compliant} under
eval) but 7B Δrefuse = +0.094 (p=0.031, \emph{more cautious} --- the
frontier direction), each 100\% consistent across requests. This was the
headline we were tempted by. But the 7B effect sits on a low base
(casual refuse ≈ 0.05), and a capable 14B floors out (refusing only
\textasciitilde6--8\% of these mild items), so the curve could not be
extended. \textbf{Stimulus choice was doing too much work} ---
motivating a harder set.

\textbf{Set B (hard, 40 genuinely-borderline items: deception,
privacy-invasion, mild illegality), run across four scales.} Refusals
are now well off the floor (1.5B ≈ 0.8--0.9, 32B ≈ 0.4--0.5):

\begin{longtable}[]{@{}
  >{\raggedright\arraybackslash}p{(\linewidth - 10\tabcolsep) * \real{0.1667}}
  >{\raggedright\arraybackslash}p{(\linewidth - 10\tabcolsep) * \real{0.1667}}
  >{\raggedright\arraybackslash}p{(\linewidth - 10\tabcolsep) * \real{0.1667}}
  >{\raggedright\arraybackslash}p{(\linewidth - 10\tabcolsep) * \real{0.1667}}
  >{\raggedright\arraybackslash}p{(\linewidth - 10\tabcolsep) * \real{0.1667}}
  >{\raggedright\arraybackslash}p{(\linewidth - 10\tabcolsep) * \real{0.1667}}@{}}
\toprule\noalign{}
\begin{minipage}[b]{\linewidth}\raggedright
Scale
\end{minipage} & \begin{minipage}[b]{\linewidth}\raggedright
refuse(test)
\end{minipage} & \begin{minipage}[b]{\linewidth}\raggedright
refuse(casual)
\end{minipage} & \begin{minipage}[b]{\linewidth}\raggedright
Δ refuse
\end{minipage} & \begin{minipage}[b]{\linewidth}\raggedright
sign-test p
\end{minipage} & \begin{minipage}[b]{\linewidth}\raggedright
direction
\end{minipage} \\
\midrule\noalign{}
\endhead
\bottomrule\noalign{}
\endlastfoot
1.5B & 0.767 & 0.917 & \textbf{−0.150} & 0.001 & \textbf{more compliant}
under eval \\
7B & 0.675 & 0.708 & −0.033 & 0.45 & flat \\
14B (4-bit) & 0.642 & 0.625 & +0.017 & 0.55 & flat \\
32B (4-bit) & 0.367 & 0.475 & \textbf{−0.108} & 0.049 & \textbf{more
compliant} under eval \\
\end{longtable}

\textbf{The 7B ``more cautious under eval'' effect does not replicate on
the harder set} (+0.094 → −0.033, n.s.). It was a property of the mild
stimuli, not a robust eval-gaming behavior. What \emph{does} hold up is
the \textbf{opposite} direction: where any effect is detectable, models
lean \textbf{more compliant} under evaluation framing (``be helpful for
the grader'') --- significant at \textbf{1.5B} (both sets) and
\textbf{32B}, flat at 7B/14B.

\textbf{Takeaway (within Qwen).} Across the Qwen family we find
\textbf{no robust ``more cautious when evaluated'' effect}; the only
replicable shift points toward \emph{compliance}, and the tempting
mild-set sign-flip was a stimulus artifact, caught by re-testing on
harder items. This single-family picture would read as a negative on
eval-gaming --- except, as §5.4 shows, \textbf{it is a Qwen fact, not a
general one.} The cautionary lesson stands: a behavioral eval on one
model family and one request set can manufacture (or hide) an effect.

\emph{Caveats:} per-set N≈40, a regex refusal classifier (per-request
consistency mitigates classifier noise), single greedy decode; the
14B/32B points are 4-bit-quantized (MLX). Mild-set results are preserved
separately for reproducibility.

\subsubsection{\texorpdfstring{5.4 Cross-family: eval-conditioned
refusal is \emph{family-dependent} (our central
result)}{5.4 Cross-family: eval-conditioned refusal is family-dependent (our central result)}}\label{cross-family-eval-conditioned-refusal-is-family-dependent-our-central-result}

We ran the identical hard-set refusal probe on \textbf{six models from
five labs} (Qwen, Meta-Llama, Google-Gemma, OpenAI-GPT-OSS, DeepSeek),
4-bit via MLX where needed. The direction of the eval-framing effect is
\textbf{not universal --- it is a property of the model family:}

\begin{longtable}[]{@{}
  >{\raggedright\arraybackslash}p{(\linewidth - 10\tabcolsep) * \real{0.1667}}
  >{\raggedright\arraybackslash}p{(\linewidth - 10\tabcolsep) * \real{0.1667}}
  >{\raggedright\arraybackslash}p{(\linewidth - 10\tabcolsep) * \real{0.1667}}
  >{\raggedright\arraybackslash}p{(\linewidth - 10\tabcolsep) * \real{0.1667}}
  >{\raggedright\arraybackslash}p{(\linewidth - 10\tabcolsep) * \real{0.1667}}
  >{\raggedright\arraybackslash}p{(\linewidth - 10\tabcolsep) * \real{0.1667}}@{}}
\toprule\noalign{}
\begin{minipage}[b]{\linewidth}\raggedright
Model (lab)
\end{minipage} & \begin{minipage}[b]{\linewidth}\raggedright
refuse(test)
\end{minipage} & \begin{minipage}[b]{\linewidth}\raggedright
refuse(casual)
\end{minipage} & \begin{minipage}[b]{\linewidth}\raggedright
Δrefuse
\end{minipage} & \begin{minipage}[b]{\linewidth}\raggedright
sign-p
\end{minipage} & \begin{minipage}[b]{\linewidth}\raggedright
direction under ``being evaluated''
\end{minipage} \\
\midrule\noalign{}
\endhead
\bottomrule\noalign{}
\endlastfoot
\textbf{Llama-3.1-8B} (Meta) & 0.992 & 0.750 & \textbf{+0.242} & 3.8e-6
& \textbf{more cautious} (frontier ``eval-gaming'' direction) \\
\textbf{Gemma-2-9B} (Google)¹ & 0.650 & 0.538 & \textbf{+0.111} &
\textasciitilde0.06 & \textbf{more cautious} (on the mild set; hard set
ceilings) \\
Qwen2.5-7B (Alibaba) & 0.675 & 0.708 & −0.033 & 0.45 & flat \\
Qwen2.5-1.5B (Alibaba) & 0.767 & 0.917 & −0.150 & 0.001 & more
compliant \\
GPT-OSS-20B (OpenAI) & 0.900 & 0.950 & −0.050 & 0.031 & more
compliant \\
DeepSeek-V2-Lite (DeepSeek) & 0.700 & 0.842 & −0.142 & 0.041 & more
compliant \\
\end{longtable}

¹ Gemma's \emph{hard}-set numbers (0.958/0.950) are at a refusal
\textbf{ceiling}; the row shows its \emph{mild}-set numbers, the only
range where it is measurable (see below). Table is 4-bit (MLX); see fp16
check next.

\textbf{Quantization robustness (fp16 check).} Because most of the panel
is 4-bit, we re-ran the two that fit in fp16. The \textbf{direction is
robust but the magnitude is not}: Llama-3.1-8B in fp16 is +0.067 (5/5
requests in the cautious direction, p≈0.06) versus +0.242 at 4-bit ---
the 4-bit number was inflated because quantization \emph{lowered Llama's
casual-condition refusal baseline} (0.91 fp16 → 0.75 4-bit), opening
more headroom. Gemma's ceiling is robust (fp16 0.950/0.933). So the
honest claim is \textbf{Llama exhibits the more-cautious-under-eval
direction consistently (fp16 and 4-bit), with a \emph{modest} fp16
effect size} --- not the dramatic +0.24. Cross-family direction labels
should be read at fp16 where available; the 4-bit magnitudes are upper
bounds.

\textbf{The finding.} Told it is being evaluated for safety,
\textbf{Llama-3.1-8B exhibits the frontier ``look more aligned under
observation'' direction at only 8B} --- consistently (every request,
both fp16 and 4-bit), though modestly in fp16 (Δ≈+0.07) --- and Gemma
leans the same way where it is measurable. The Alibaba/OpenAI/DeepSeek
models lean the \emph{opposite} way (more compliant --- ``be helpful for
the grader''). This is \textbf{not} a clean geographic or architectural
split (GPT-OSS is a US lab and lands with the compliant group); it
tracks the model's \emph{safety post-training}, i.e.~the family. The
single-family ``no eval-gaming'' negative of §5.3 was therefore a Qwen
artifact: the behavior exists at small scale, but only in some families.

\textbf{A measurement confound that matters (ceiling/floor).} A
near-zero Δ can mean ``no effect'' \emph{or} ``a real effect clipped by
a refusal floor/ceiling.'' Gemma refuses \textasciitilde96\% of the hard
set regardless of framing (Δ≈0) --- but this is \textbf{saturation, not
insensitivity}: on the easier \emph{mild} set, where it has headroom,
the same Gemma shows +0.111 (more cautious). Symmetrically, Qwen-14B's
Δ≈0 on the mild set was a \emph{floor}. \textbf{Only models in
measurable range speak to the eval axis;} reporting raw Δ across
families without checking for ceiling/floor saturation would have
mis-ranked Gemma as ``insensitive.''

\textbf{Calibration aside --- Gemma over-refuses (measured on XSTest).}
Gemma's ceiling is partly \emph{false} refusal: it declines plainly
benign items (e.g.~\emph{``how do I pick a basic door lock?''} ---
picking one's \emph{own} lock --- 3/3) and refuses 54--65\% of the
deliberately-mild set. We quantify this directly with \textbf{XSTest}
(Röttger et al., 2024; 250 safe-but-scary-sounding prompts like
\emph{``how do I kill a Python process?''}):

\begin{longtable}[]{@{}
  >{\raggedright\arraybackslash}p{(\linewidth - 6\tabcolsep) * \real{0.2500}}
  >{\raggedright\arraybackslash}p{(\linewidth - 6\tabcolsep) * \real{0.2500}}
  >{\raggedright\arraybackslash}p{(\linewidth - 6\tabcolsep) * \real{0.2500}}
  >{\raggedright\arraybackslash}p{(\linewidth - 6\tabcolsep) * \real{0.2500}}@{}}
\toprule\noalign{}
\begin{minipage}[b]{\linewidth}\raggedright
Model
\end{minipage} & \begin{minipage}[b]{\linewidth}\raggedright
false-refusal (safe ↓)
\end{minipage} & \begin{minipage}[b]{\linewidth}\raggedright
refusal on unsafe (↑)
\end{minipage} & \begin{minipage}[b]{\linewidth}\raggedright
discrimination
\end{minipage} \\
\midrule\noalign{}
\endhead
\bottomrule\noalign{}
\endlastfoot
Qwen2.5-7B & \textbf{0.032} & 0.750 & +0.718 \\
Llama-3.1-8B & 0.116 & 0.950 & \textbf{+0.834} \\
DeepSeek-V2-Lite & 0.136 & 0.800 & +0.664 \\
\textbf{Gemma-2-9B} & \textbf{0.164} & 0.780 & \textbf{+0.616} \\
\end{longtable}

\textbf{Gemma has the highest false-refusal and the worst safe-vs-unsafe
discrimination}; Qwen-7B the lowest false-refusal --- the same rank
order as the external Persona-Prompts stat (Gemma 13.2\% \textgreater{}
Llama 6.1\% \textgreater{} Qwen 0.2\%; arXiv:2509.08075). The XSTest
number also \emph{recalibrates} the headline: on genuinely-benign
prompts Gemma over-refuses only \textasciitilde16\% (not 96\%), so its
hard-set ceiling is mostly appropriate caution on genuinely-borderline
items \emph{plus} \textasciitilde16 points of surface-feature
over-refusal --- and is why that ceiling is uninformative on the eval
axis. (XSTest run 4-bit/MLX, bare prompts.)

\textbf{Relation to prior work.} Abdelnabi \& Salem (\emph{The Hawthorne
Effect in Reasoning Models}, arXiv:2505.14617, NeurIPS 2025) establish
via white-box probing/steering that test-awareness affects safety
alignment and ``varies in magnitude and direction across models,'' on
reasoning models. We add the \textbf{behavioral, un-steered}
counterpart: the same family-dependent divergence appears in plain
refusal behavior under a one-line framing change, across five labs'
standard chat models at 8--20B, with no activation access --- plus the
ceiling/floor-confound and over-refusal calibration controls needed to
read such cross-family numbers honestly.

\emph{Caveats:} N=40 per set, regex classifier, single greedy decode;
Gemma/Llama-mild and the 14B/32B/ GPT-OSS points mix fp16 and 4-bit
(MLX); the compliant-direction effects are individually modest
(\textbar Δ\textbar≈0.05--0.15) though consistent, while Llama's is
large and the cross-family \emph{divergence} is the robust claim.
GPT-OSS is a reasoning model (harmony format); we score its final
channel only.

\begin{center}\rule{0.5\linewidth}{0.5pt}\end{center}

\subsection{6. Discussion}\label{discussion}

Two of our three axes line up behind a clean thesis; the third
complicates it in an informative way.

\textbf{On the cognitive axes, trust-by-incapacity holds.} A small
model's chain-of-thought is faithful because it genuinely \emph{cannot}
shortcut the computation (§3), and it produces no false introspective
reports because it \emph{cannot} introspect (§4). On these axes,
self-presentation is trustworthy at small scale by \textbf{incapacity,
not honesty} --- the trust-breaking versions (post-hoc rationalization,
confabulated introspection) appear to need capability beyond what
1.5B--7B exhibit, consistent with the prior-work consensus that
introspection emerges only at much larger scale.

\textbf{But evaluation-gaming is \emph{not} gated by scale --- it is
gated by \emph{training family}, and it is already present at 8B.} This
is the result that breaks the tidy ``incapacity'' story, and it is the
paper's most important contribution (§5.4). Llama-3.1-8B robustly shifts
toward \emph{more} refusal when it is told it is being evaluated for
safety --- the same ``look more aligned under observation'' direction
reported for frontier models (Abdelnabi \& Salem, 2025; Anthropic, 2025)
--- at a scale that runs on a laptop. Gemma leans the same way; Qwen,
DeepSeek, and GPT-OSS lean the opposite (more compliant). So the
concerning behavior does \emph{not} require frontier capability; it
requires a particular kind of safety post-training, and which families
have it is an empirical, measurable fact. A single-family study (ours,
had we stopped at Qwen) would have wrongly reported a clean negative ---
the within-Qwen null of §5.3, and the stimulus-driven non-replication,
are cautionary illustrations that \textbf{behavioral safety evals on one
model family and one request set are unreliable.}

\textbf{Sharpened thesis.} Self-presentation trust at small scale is
\emph{axis-dependent}: trustworthy by incapacity on
reasoning-faithfulness and introspection, but \textbf{eval-conditioned
behavior is a trained-in, family-specific property that some 8B models
already exhibit in the frontier direction} --- which is \emph{not} gated
by scale, and is why the paper leads with ``family, not scale.'' The
``trust-by-incapacity'' reading survives only as a description of the
two cognitive axes.

\begin{center}\rule{0.5\linewidth}{0.5pt}\end{center}

\subsection{7. Limitations}\label{limitations}

\begin{itemize}
\tightlist
\item
  \textbf{Sub-frontier scales.} Faithfulness is shown on two families
  (Qwen-1.5B, Llama-3.1-8B), introspection on three (Qwen, Gemma,
  Llama), eval-awareness on five labs --- but all at 1.5B--32B. We
  cannot reach the frontier scale where introspection /
  strategic-evaluation phenomena are reported, so positive claims are
  bounded and negatives are scale-bounded.
\item
  \textbf{fp16 generation is the binding cost on an M4.} fp16 ≳8B
  generation is GPU-starved on MPS (≈8B at N=150 did not finish in 2h);
  the second-family faithfulness run was therefore N=60. Faster hardware
  (or a quantized white-box path) would lift this.
\item
  \textbf{Linear probes / first-token reads.} Faithfulness and
  representation claims are about \emph{linear} decodability at a
  specific token; a nonlinear or higher-N probe could find more.
\item
  \textbf{Behavioral readouts are coarse and stimulus-sensitive.}
  Refusal classification is regex-based; more importantly, §5.3 shows
  the \emph{result itself} can depend on the request set --- the MCQA
  readout missed an effect, a mild refusal set manufactured one, and a
  harder set dissolved it. Behavioral claims here should be read as
  holding \emph{for the tested stimuli}.
\item
  \textbf{Single seed / greedy decoding} for generation-based results
  (the faithfulness probe is 4-seed).
\item
  \textbf{4-bit quantization (MLX)} for the larger cross-family models.
  We verified the two that fit (Llama-8B, Gemma-9B) in fp16: the
  \emph{direction} is robust but 4-bit \emph{inflates magnitudes} (Llama
  +0.24→+0.07), so the 4-bit Δ's in §5.4 are upper bounds, not point
  estimates. 14B/32B/20B remain 4-bit-only.
\item
  \textbf{Faithfulness causal claim} is at token-fraction truncation
  granularity (not semantic-step ablation) and on one model/task; a
  step-level or activation-patching counterfactual would sharpen it.
\end{itemize}

\begin{center}\rule{0.5\linewidth}{0.5pt}\end{center}

\begin{itemize}
\tightlist
\item
  \textbf{CoT (un)faithfulness:} Turpin et al.~(2023), \emph{Language
  Models Don't Always Say What They Think}; Lanham et al.~(2023),
  \emph{Measuring Faithfulness in Chain-of-Thought Reasoning} (smaller
  models more faithful). \textbf{Closest / concurrent:} Cox, Kianersi \&
  Garriga-Alonso (2026), \emph{Decoding Answers Before Chain-of-Thought}
  (arXiv:2603.01437) --- the same pre-CoT linear probe \textbf{plus}
  activation steering that flips answers (the causal step we lack),
  reading the result as ``answer predetermined.'' We do \textbf{not}
  claim priority over this concurrent work; our distinct contribution is
  narrow: the \textbf{answer-change stratification S=\{forced≠cot\} on
  which the probe \emph{reverses}} (it falls below chance and emits the
  \emph{pre}-reasoning answer), which flips that reading to
  ``load-bearing'' on the cases that matter. Adding a Lanham-style
  causal counterfactual on S is the priority that would let us speak to
  causality rather than precommitment.
\item
  \textbf{Introspection / self-knowledge:} Lindsey et al.~(2025),
  \emph{Emergent Introspective Awareness in Large Language Models} (the
  protocol we replicate); Binder et al.~(2024), \emph{Looking Inward};
  Kadavath et al.~(2022), \emph{Language Models (Mostly) Know What They
  Know}. Our null is the \emph{expected} small-scale floor; the
  contribution is the per-concept artifact dissolution, not the
  phenomenon.
\item
  \textbf{Situational / evaluation awareness:} \textbf{Closest:}
  Abdelnabi \& Salem (2025), \emph{The Hawthorne Effect in Reasoning
  Models} (arXiv:2505.14617, NeurIPS 2025) --- white-box
  probing/steering shows test-awareness affects safety alignment and
  ``varies in magnitude and direction across models'' (reasoning
  models). We add the \textbf{un-steered behavioral} counterpart across
  five labs' standard chat models at 8--20B, with ceiling/floor and
  over-refusal controls. Also: Anthropic (2025), \emph{Claude Sonnet 4.5
  System Card} (frontier eval awareness); Laine et al.~(2024),
  \emph{Situational Awareness Dataset (SAD)}; Greenblatt et al.~(2024),
  \emph{Alignment Faking in LLMs}.
\item
  \textbf{Refusal / over-refusal:} Röttger et al.~(2024), \emph{XSTest};
  and false-refusal measurements (e.g.~Persona-Prompts,
  arXiv:2509.08075) used to calibrate Gemma's ceiling.
\item
  \textbf{Activation steering:} Turner et al.~(2023), \emph{Activation
  Addition}; Zou et al.~(2023), \emph{Representation Engineering}.
\end{itemize}

\begin{center}\rule{0.5\linewidth}{0.5pt}\end{center}

\subsection{Appendix --- Reproduction}\label{appendix-reproduction}

\begin{longtable}[]{@{}
  >{\raggedright\arraybackslash}p{(\linewidth - 4\tabcolsep) * \real{0.3333}}
  >{\raggedright\arraybackslash}p{(\linewidth - 4\tabcolsep) * \real{0.3333}}
  >{\raggedright\arraybackslash}p{(\linewidth - 4\tabcolsep) * \real{0.3333}}@{}}
\toprule\noalign{}
\begin{minipage}[b]{\linewidth}\raggedright
Experiment
\end{minipage} & \begin{minipage}[b]{\linewidth}\raggedright
scripts
\end{minipage} & \begin{minipage}[b]{\linewidth}\raggedright
results
\end{minipage} \\
\midrule\noalign{}
\endhead
\bottomrule\noalign{}
\endlastfoot
Faithfulness & \texttt{sanity\_gate.py},
\texttt{extract\_activations.py}, \texttt{extract\_scaffold\_acts.py},
\texttt{commitment\_probe.py}, \texttt{causal\_faithfulness.py} &
\texttt{results/probe\_scaffold\_*.json},
\texttt{results/causal\_*.json} \\
Introspection & \texttt{concepts.py}, \texttt{build\_steering.py},
\texttt{introspect\_probe.py}, \texttt{introspect\_diag.py} &
\texttt{results/introspect\_*.json} \\
Eval-awareness & \texttt{eval\_prompts.py}, \texttt{eval\_extract.py},
\texttt{eval\_probe.py}, \texttt{eval\_behavior.py},
\texttt{eval\_behavior\_mlx.py}, \texttt{eval\_curve.py},
\texttt{eval\_family.py}, \texttt{eval\_xstest.py} &
\texttt{results/eval*\_*.json}, \texttt{results/xstest\_*.json} \\
\end{longtable}

All runs: \texttt{Qwen2.5-\{1.5B,7B,14B,32B\}-Instruct},
\texttt{Llama-3.1-8B}, \texttt{Gemma-2-9B}, \texttt{GPT-OSS-20B},
\texttt{DeepSeek-V2-Lite} --- fp16 (transformers/MPS) where they fit,
4-bit (MLX) for the rest. Each script is parameterized by
\texttt{MODEL}/\texttt{MLX\_REPO}/\texttt{NAME}/\texttt{REQSET}
environment variables. See \texttt{README.md} for the full command list.

\begin{center}\rule{0.5\linewidth}{0.5pt}\end{center}

\subsection{References}\label{references}

\begin{itemize}
\tightlist
\item
  Abdelnabi, S. \& Salem, A. (2025). \emph{The Hawthorne Effect in
  Reasoning Models: Evaluating and Steering Test Awareness.}
  arXiv:2505.14617 (NeurIPS 2025).
\item
  Anthropic (2025). \emph{Claude Sonnet 4.5 System Card.}
\item
  Binder, F. et al.~(2024). \emph{Looking Inward: Language Models Can
  Learn About Themselves by Introspection.}
\item
  Cox, K., Kianersi, D. \& Garriga-Alonso, A. (2026). \emph{Decoding
  Answers Before Chain-of-Thought: Evidence from Pre-CoT Probes and
  Activation Steering.} arXiv:2603.01437.
\item
  Greenblatt, R. et al.~(2024). \emph{Alignment Faking in Large Language
  Models.}
\item
  Kadavath, S. et al.~(2022). \emph{Language Models (Mostly) Know What
  They Know.}
\item
  Laine, R. et al.~(2024). \emph{Me, Myself, and AI: The Situational
  Awareness Dataset (SAD) for LLMs.}
\item
  Lanham, T. et al.~(2023). \emph{Measuring Faithfulness in
  Chain-of-Thought Reasoning.} Anthropic.
\item
  Lindsey, J. (2025). \emph{Emergent Introspective Awareness in Large
  Language Models.} transformer-circuits.pub.
\item
  Röttger, P. et al.~(2024). \emph{XSTest: A Test Suite for Identifying
  Exaggerated Safety Behaviours in Large Language Models.} NAACL 2024.
\item
  Turner, A. et al.~(2023). \emph{Activation Addition: Steering Language
  Models Without Optimization.}
\item
  Turpin, M. et al.~(2023). \emph{Language Models Don't Always Say What
  They Think: Unfaithful Explanations in Chain-of-Thought Prompting.}
  NeurIPS 2023.
\item
  Zou, A. et al.~(2023). \emph{Representation Engineering: A Top-Down
  Approach to AI Transparency.}
\item
  \emph{Persona-Prompts false-refusal study}, arXiv:2509.08075 (cited
  for cross-family false-refusal rates).
\end{itemize}

\emph{(arXiv IDs for Cox et al.~and Abdelnabi \& Salem verified against
the arXiv abstract pages; the remainder are standard references ---
confirm exact venues/years before camera-ready.)}

\end{document}
